\documentclass[11pt]{article}

    \usepackage[breakable]{tcolorbox}
    \usepackage{parskip} % Stop auto-indenting (to mimic markdown behaviour)
    

    % Basic figure setup, for now with no caption control since it's done
    % automatically by Pandoc (which extracts ![](path) syntax from Markdown).
    \usepackage{graphicx}
    % Keep aspect ratio if custom image width or height is specified
    \setkeys{Gin}{keepaspectratio}
    % Maintain compatibility with old templates. Remove in nbconvert 6.0
    \let\Oldincludegraphics\includegraphics
    % Ensure that by default, figures have no caption (until we provide a
    % proper Figure object with a Caption API and a way to capture that
    % in the conversion process - todo).
    \usepackage{caption}
    \DeclareCaptionFormat{nocaption}{}
    \captionsetup{format=nocaption,aboveskip=0pt,belowskip=0pt}

    \usepackage{float}
    \floatplacement{figure}{H} % forces figures to be placed at the correct location
    \usepackage{xcolor} % Allow colors to be defined
    \usepackage{enumerate} % Needed for markdown enumerations to work
    \usepackage{geometry} % Used to adjust the document margins
    \usepackage{amsmath} % Equations
    \usepackage{amssymb} % Equations
    \usepackage{textcomp} % defines textquotesingle
    % Hack from http://tex.stackexchange.com/a/47451/13684:
    \AtBeginDocument{%
        \def\PYZsq{\textquotesingle}% Upright quotes in Pygmentized code
    }
    \usepackage{upquote} % Upright quotes for verbatim code
    \usepackage{eurosym} % defines \euro

    \usepackage{iftex}
    \ifPDFTeX
        \usepackage[T1]{fontenc}
        \IfFileExists{alphabeta.sty}{
              \usepackage{alphabeta}
          }{
              \usepackage[mathletters]{ucs}
              \usepackage[utf8x]{inputenc}
          }
    \else
        \usepackage{fontspec}
        \usepackage{unicode-math}
    \fi

    \usepackage{fancyvrb} % verbatim replacement that allows latex
    \usepackage{grffile} % extends the file name processing of package graphics
                         % to support a larger range
    \makeatletter % fix for old versions of grffile with XeLaTeX
    \@ifpackagelater{grffile}{2019/11/01}
    {
      % Do nothing on new versions
    }
    {
      \def\Gread@@xetex#1{%
        \IfFileExists{"\Gin@base".bb}%
        {\Gread@eps{\Gin@base.bb}}%
        {\Gread@@xetex@aux#1}%
      }
    }
    \makeatother
    \usepackage[Export]{adjustbox} % Used to constrain images to a maximum size
    \adjustboxset{max size={0.9\linewidth}{0.9\paperheight}}

    % The hyperref package gives us a pdf with properly built
    % internal navigation ('pdf bookmarks' for the table of contents,
    % internal cross-reference links, web links for URLs, etc.)
    \usepackage{hyperref}
    % The default LaTeX title has an obnoxious amount of whitespace. By default,
    % titling removes some of it. It also provides customization options.
    \usepackage{titling}
    \usepackage{longtable} % longtable support required by pandoc >1.10
    \usepackage{booktabs}  % table support for pandoc > 1.12.2
    \usepackage{array}     % table support for pandoc >= 2.11.3
    \usepackage{calc}      % table minipage width calculation for pandoc >= 2.11.1
    \usepackage[inline]{enumitem} % IRkernel/repr support (it uses the enumerate* environment)
    \usepackage[normalem]{ulem} % ulem is needed to support strikethroughs (\sout)
                                % normalem makes italics be italics, not underlines
    \usepackage{soul}      % strikethrough (\st) support for pandoc >= 3.0.0
    \usepackage{mathrsfs}
    

    
    % Colors for the hyperref package
    \definecolor{urlcolor}{rgb}{0,.145,.698}
    \definecolor{linkcolor}{rgb}{.71,0.21,0.01}
    \definecolor{citecolor}{rgb}{.12,.54,.11}

    % ANSI colors
    \definecolor{ansi-black}{HTML}{3E424D}
    \definecolor{ansi-black-intense}{HTML}{282C36}
    \definecolor{ansi-red}{HTML}{E75C58}
    \definecolor{ansi-red-intense}{HTML}{B22B31}
    \definecolor{ansi-green}{HTML}{00A250}
    \definecolor{ansi-green-intense}{HTML}{007427}
    \definecolor{ansi-yellow}{HTML}{DDB62B}
    \definecolor{ansi-yellow-intense}{HTML}{B27D12}
    \definecolor{ansi-blue}{HTML}{208FFB}
    \definecolor{ansi-blue-intense}{HTML}{0065CA}
    \definecolor{ansi-magenta}{HTML}{D160C4}
    \definecolor{ansi-magenta-intense}{HTML}{A03196}
    \definecolor{ansi-cyan}{HTML}{60C6C8}
    \definecolor{ansi-cyan-intense}{HTML}{258F8F}
    \definecolor{ansi-white}{HTML}{C5C1B4}
    \definecolor{ansi-white-intense}{HTML}{A1A6B2}
    \definecolor{ansi-default-inverse-fg}{HTML}{FFFFFF}
    \definecolor{ansi-default-inverse-bg}{HTML}{000000}

    % common color for the border for error outputs.
    \definecolor{outerrorbackground}{HTML}{FFDFDF}

    % commands and environments needed by pandoc snippets
    % extracted from the output of `pandoc -s`
    \providecommand{\tightlist}{%
      \setlength{\itemsep}{0pt}\setlength{\parskip}{0pt}}
    \DefineVerbatimEnvironment{Highlighting}{Verbatim}{commandchars=\\\{\}}
    % Add ',fontsize=\small' for more characters per line
    \newenvironment{Shaded}{}{}
    \newcommand{\KeywordTok}[1]{\textcolor[rgb]{0.00,0.44,0.13}{\textbf{{#1}}}}
    \newcommand{\DataTypeTok}[1]{\textcolor[rgb]{0.56,0.13,0.00}{{#1}}}
    \newcommand{\DecValTok}[1]{\textcolor[rgb]{0.25,0.63,0.44}{{#1}}}
    \newcommand{\BaseNTok}[1]{\textcolor[rgb]{0.25,0.63,0.44}{{#1}}}
    \newcommand{\FloatTok}[1]{\textcolor[rgb]{0.25,0.63,0.44}{{#1}}}
    \newcommand{\CharTok}[1]{\textcolor[rgb]{0.25,0.44,0.63}{{#1}}}
    \newcommand{\StringTok}[1]{\textcolor[rgb]{0.25,0.44,0.63}{{#1}}}
    \newcommand{\CommentTok}[1]{\textcolor[rgb]{0.38,0.63,0.69}{\textit{{#1}}}}
    \newcommand{\OtherTok}[1]{\textcolor[rgb]{0.00,0.44,0.13}{{#1}}}
    \newcommand{\AlertTok}[1]{\textcolor[rgb]{1.00,0.00,0.00}{\textbf{{#1}}}}
    \newcommand{\FunctionTok}[1]{\textcolor[rgb]{0.02,0.16,0.49}{{#1}}}
    \newcommand{\RegionMarkerTok}[1]{{#1}}
    \newcommand{\ErrorTok}[1]{\textcolor[rgb]{1.00,0.00,0.00}{\textbf{{#1}}}}
    \newcommand{\NormalTok}[1]{{#1}}

    % Additional commands for more recent versions of Pandoc
    \newcommand{\ConstantTok}[1]{\textcolor[rgb]{0.53,0.00,0.00}{{#1}}}
    \newcommand{\SpecialCharTok}[1]{\textcolor[rgb]{0.25,0.44,0.63}{{#1}}}
    \newcommand{\VerbatimStringTok}[1]{\textcolor[rgb]{0.25,0.44,0.63}{{#1}}}
    \newcommand{\SpecialStringTok}[1]{\textcolor[rgb]{0.73,0.40,0.53}{{#1}}}
    \newcommand{\ImportTok}[1]{{#1}}
    \newcommand{\DocumentationTok}[1]{\textcolor[rgb]{0.73,0.13,0.13}{\textit{{#1}}}}
    \newcommand{\AnnotationTok}[1]{\textcolor[rgb]{0.38,0.63,0.69}{\textbf{\textit{{#1}}}}}
    \newcommand{\CommentVarTok}[1]{\textcolor[rgb]{0.38,0.63,0.69}{\textbf{\textit{{#1}}}}}
    \newcommand{\VariableTok}[1]{\textcolor[rgb]{0.10,0.09,0.49}{{#1}}}
    \newcommand{\ControlFlowTok}[1]{\textcolor[rgb]{0.00,0.44,0.13}{\textbf{{#1}}}}
    \newcommand{\OperatorTok}[1]{\textcolor[rgb]{0.40,0.40,0.40}{{#1}}}
    \newcommand{\BuiltInTok}[1]{{#1}}
    \newcommand{\ExtensionTok}[1]{{#1}}
    \newcommand{\PreprocessorTok}[1]{\textcolor[rgb]{0.74,0.48,0.00}{{#1}}}
    \newcommand{\AttributeTok}[1]{\textcolor[rgb]{0.49,0.56,0.16}{{#1}}}
    \newcommand{\InformationTok}[1]{\textcolor[rgb]{0.38,0.63,0.69}{\textbf{\textit{{#1}}}}}
    \newcommand{\WarningTok}[1]{\textcolor[rgb]{0.38,0.63,0.69}{\textbf{\textit{{#1}}}}}
    \makeatletter
    \newsavebox\pandoc@box
    \newcommand*\pandocbounded[1]{%
      \sbox\pandoc@box{#1}%
      % scaling factors for width and height
      \Gscale@div\@tempa\textheight{\dimexpr\ht\pandoc@box+\dp\pandoc@box\relax}%
      \Gscale@div\@tempb\linewidth{\wd\pandoc@box}%
      % select the smaller of both
      \ifdim\@tempb\p@<\@tempa\p@
        \let\@tempa\@tempb
      \fi
      % scaling accordingly (\@tempa < 1)
      \ifdim\@tempa\p@<\p@
        \scalebox{\@tempa}{\usebox\pandoc@box}%
      % scaling not needed, use as it is
      \else
        \usebox{\pandoc@box}%
      \fi
    }
    \makeatother

    % Define a nice break command that doesn't care if a line doesn't already
    % exist.
    \def\br{\hspace*{\fill} \\* }
    % Math Jax compatibility definitions
    \def\gt{>}
    \def\lt{<}
    \let\Oldtex\TeX
    \let\Oldlatex\LaTeX
    \renewcommand{\TeX}{\textrm{\Oldtex}}
    \renewcommand{\LaTeX}{\textrm{\Oldlatex}}
    % Document parameters
    % Document title
    \title{evaluate\_external}
    
    
    
    
    
    
    
% Pygments definitions
\makeatletter
\def\PY@reset{\let\PY@it=\relax \let\PY@bf=\relax%
    \let\PY@ul=\relax \let\PY@tc=\relax%
    \let\PY@bc=\relax \let\PY@ff=\relax}
\def\PY@tok#1{\csname PY@tok@#1\endcsname}
\def\PY@toks#1+{\ifx\relax#1\empty\else%
    \PY@tok{#1}\expandafter\PY@toks\fi}
\def\PY@do#1{\PY@bc{\PY@tc{\PY@ul{%
    \PY@it{\PY@bf{\PY@ff{#1}}}}}}}
\def\PY#1#2{\PY@reset\PY@toks#1+\relax+\PY@do{#2}}

\@namedef{PY@tok@w}{\def\PY@tc##1{\textcolor[rgb]{0.73,0.73,0.73}{##1}}}
\@namedef{PY@tok@c}{\let\PY@it=\textit\def\PY@tc##1{\textcolor[rgb]{0.24,0.48,0.48}{##1}}}
\@namedef{PY@tok@cp}{\def\PY@tc##1{\textcolor[rgb]{0.61,0.40,0.00}{##1}}}
\@namedef{PY@tok@k}{\let\PY@bf=\textbf\def\PY@tc##1{\textcolor[rgb]{0.00,0.50,0.00}{##1}}}
\@namedef{PY@tok@kp}{\def\PY@tc##1{\textcolor[rgb]{0.00,0.50,0.00}{##1}}}
\@namedef{PY@tok@kt}{\def\PY@tc##1{\textcolor[rgb]{0.69,0.00,0.25}{##1}}}
\@namedef{PY@tok@o}{\def\PY@tc##1{\textcolor[rgb]{0.40,0.40,0.40}{##1}}}
\@namedef{PY@tok@ow}{\let\PY@bf=\textbf\def\PY@tc##1{\textcolor[rgb]{0.67,0.13,1.00}{##1}}}
\@namedef{PY@tok@nb}{\def\PY@tc##1{\textcolor[rgb]{0.00,0.50,0.00}{##1}}}
\@namedef{PY@tok@nf}{\def\PY@tc##1{\textcolor[rgb]{0.00,0.00,1.00}{##1}}}
\@namedef{PY@tok@nc}{\let\PY@bf=\textbf\def\PY@tc##1{\textcolor[rgb]{0.00,0.00,1.00}{##1}}}
\@namedef{PY@tok@nn}{\let\PY@bf=\textbf\def\PY@tc##1{\textcolor[rgb]{0.00,0.00,1.00}{##1}}}
\@namedef{PY@tok@ne}{\let\PY@bf=\textbf\def\PY@tc##1{\textcolor[rgb]{0.80,0.25,0.22}{##1}}}
\@namedef{PY@tok@nv}{\def\PY@tc##1{\textcolor[rgb]{0.10,0.09,0.49}{##1}}}
\@namedef{PY@tok@no}{\def\PY@tc##1{\textcolor[rgb]{0.53,0.00,0.00}{##1}}}
\@namedef{PY@tok@nl}{\def\PY@tc##1{\textcolor[rgb]{0.46,0.46,0.00}{##1}}}
\@namedef{PY@tok@ni}{\let\PY@bf=\textbf\def\PY@tc##1{\textcolor[rgb]{0.44,0.44,0.44}{##1}}}
\@namedef{PY@tok@na}{\def\PY@tc##1{\textcolor[rgb]{0.41,0.47,0.13}{##1}}}
\@namedef{PY@tok@nt}{\let\PY@bf=\textbf\def\PY@tc##1{\textcolor[rgb]{0.00,0.50,0.00}{##1}}}
\@namedef{PY@tok@nd}{\def\PY@tc##1{\textcolor[rgb]{0.67,0.13,1.00}{##1}}}
\@namedef{PY@tok@s}{\def\PY@tc##1{\textcolor[rgb]{0.73,0.13,0.13}{##1}}}
\@namedef{PY@tok@sd}{\let\PY@it=\textit\def\PY@tc##1{\textcolor[rgb]{0.73,0.13,0.13}{##1}}}
\@namedef{PY@tok@si}{\let\PY@bf=\textbf\def\PY@tc##1{\textcolor[rgb]{0.64,0.35,0.47}{##1}}}
\@namedef{PY@tok@se}{\let\PY@bf=\textbf\def\PY@tc##1{\textcolor[rgb]{0.67,0.36,0.12}{##1}}}
\@namedef{PY@tok@sr}{\def\PY@tc##1{\textcolor[rgb]{0.64,0.35,0.47}{##1}}}
\@namedef{PY@tok@ss}{\def\PY@tc##1{\textcolor[rgb]{0.10,0.09,0.49}{##1}}}
\@namedef{PY@tok@sx}{\def\PY@tc##1{\textcolor[rgb]{0.00,0.50,0.00}{##1}}}
\@namedef{PY@tok@m}{\def\PY@tc##1{\textcolor[rgb]{0.40,0.40,0.40}{##1}}}
\@namedef{PY@tok@gh}{\let\PY@bf=\textbf\def\PY@tc##1{\textcolor[rgb]{0.00,0.00,0.50}{##1}}}
\@namedef{PY@tok@gu}{\let\PY@bf=\textbf\def\PY@tc##1{\textcolor[rgb]{0.50,0.00,0.50}{##1}}}
\@namedef{PY@tok@gd}{\def\PY@tc##1{\textcolor[rgb]{0.63,0.00,0.00}{##1}}}
\@namedef{PY@tok@gi}{\def\PY@tc##1{\textcolor[rgb]{0.00,0.52,0.00}{##1}}}
\@namedef{PY@tok@gr}{\def\PY@tc##1{\textcolor[rgb]{0.89,0.00,0.00}{##1}}}
\@namedef{PY@tok@ge}{\let\PY@it=\textit}
\@namedef{PY@tok@gs}{\let\PY@bf=\textbf}
\@namedef{PY@tok@ges}{\let\PY@bf=\textbf\let\PY@it=\textit}
\@namedef{PY@tok@gp}{\let\PY@bf=\textbf\def\PY@tc##1{\textcolor[rgb]{0.00,0.00,0.50}{##1}}}
\@namedef{PY@tok@go}{\def\PY@tc##1{\textcolor[rgb]{0.44,0.44,0.44}{##1}}}
\@namedef{PY@tok@gt}{\def\PY@tc##1{\textcolor[rgb]{0.00,0.27,0.87}{##1}}}
\@namedef{PY@tok@err}{\def\PY@bc##1{{\setlength{\fboxsep}{\string -\fboxrule}\fcolorbox[rgb]{1.00,0.00,0.00}{1,1,1}{\strut ##1}}}}
\@namedef{PY@tok@kc}{\let\PY@bf=\textbf\def\PY@tc##1{\textcolor[rgb]{0.00,0.50,0.00}{##1}}}
\@namedef{PY@tok@kd}{\let\PY@bf=\textbf\def\PY@tc##1{\textcolor[rgb]{0.00,0.50,0.00}{##1}}}
\@namedef{PY@tok@kn}{\let\PY@bf=\textbf\def\PY@tc##1{\textcolor[rgb]{0.00,0.50,0.00}{##1}}}
\@namedef{PY@tok@kr}{\let\PY@bf=\textbf\def\PY@tc##1{\textcolor[rgb]{0.00,0.50,0.00}{##1}}}
\@namedef{PY@tok@bp}{\def\PY@tc##1{\textcolor[rgb]{0.00,0.50,0.00}{##1}}}
\@namedef{PY@tok@fm}{\def\PY@tc##1{\textcolor[rgb]{0.00,0.00,1.00}{##1}}}
\@namedef{PY@tok@vc}{\def\PY@tc##1{\textcolor[rgb]{0.10,0.09,0.49}{##1}}}
\@namedef{PY@tok@vg}{\def\PY@tc##1{\textcolor[rgb]{0.10,0.09,0.49}{##1}}}
\@namedef{PY@tok@vi}{\def\PY@tc##1{\textcolor[rgb]{0.10,0.09,0.49}{##1}}}
\@namedef{PY@tok@vm}{\def\PY@tc##1{\textcolor[rgb]{0.10,0.09,0.49}{##1}}}
\@namedef{PY@tok@sa}{\def\PY@tc##1{\textcolor[rgb]{0.73,0.13,0.13}{##1}}}
\@namedef{PY@tok@sb}{\def\PY@tc##1{\textcolor[rgb]{0.73,0.13,0.13}{##1}}}
\@namedef{PY@tok@sc}{\def\PY@tc##1{\textcolor[rgb]{0.73,0.13,0.13}{##1}}}
\@namedef{PY@tok@dl}{\def\PY@tc##1{\textcolor[rgb]{0.73,0.13,0.13}{##1}}}
\@namedef{PY@tok@s2}{\def\PY@tc##1{\textcolor[rgb]{0.73,0.13,0.13}{##1}}}
\@namedef{PY@tok@sh}{\def\PY@tc##1{\textcolor[rgb]{0.73,0.13,0.13}{##1}}}
\@namedef{PY@tok@s1}{\def\PY@tc##1{\textcolor[rgb]{0.73,0.13,0.13}{##1}}}
\@namedef{PY@tok@mb}{\def\PY@tc##1{\textcolor[rgb]{0.40,0.40,0.40}{##1}}}
\@namedef{PY@tok@mf}{\def\PY@tc##1{\textcolor[rgb]{0.40,0.40,0.40}{##1}}}
\@namedef{PY@tok@mh}{\def\PY@tc##1{\textcolor[rgb]{0.40,0.40,0.40}{##1}}}
\@namedef{PY@tok@mi}{\def\PY@tc##1{\textcolor[rgb]{0.40,0.40,0.40}{##1}}}
\@namedef{PY@tok@il}{\def\PY@tc##1{\textcolor[rgb]{0.40,0.40,0.40}{##1}}}
\@namedef{PY@tok@mo}{\def\PY@tc##1{\textcolor[rgb]{0.40,0.40,0.40}{##1}}}
\@namedef{PY@tok@ch}{\let\PY@it=\textit\def\PY@tc##1{\textcolor[rgb]{0.24,0.48,0.48}{##1}}}
\@namedef{PY@tok@cm}{\let\PY@it=\textit\def\PY@tc##1{\textcolor[rgb]{0.24,0.48,0.48}{##1}}}
\@namedef{PY@tok@cpf}{\let\PY@it=\textit\def\PY@tc##1{\textcolor[rgb]{0.24,0.48,0.48}{##1}}}
\@namedef{PY@tok@c1}{\let\PY@it=\textit\def\PY@tc##1{\textcolor[rgb]{0.24,0.48,0.48}{##1}}}
\@namedef{PY@tok@cs}{\let\PY@it=\textit\def\PY@tc##1{\textcolor[rgb]{0.24,0.48,0.48}{##1}}}

\def\PYZbs{\char`\\}
\def\PYZus{\char`\_}
\def\PYZob{\char`\{}
\def\PYZcb{\char`\}}
\def\PYZca{\char`\^}
\def\PYZam{\char`\&}
\def\PYZlt{\char`\<}
\def\PYZgt{\char`\>}
\def\PYZsh{\char`\#}
\def\PYZpc{\char`\%}
\def\PYZdl{\char`\$}
\def\PYZhy{\char`\-}
\def\PYZsq{\char`\'}
\def\PYZdq{\char`\"}
\def\PYZti{\char`\~}
% for compatibility with earlier versions
\def\PYZat{@}
\def\PYZlb{[}
\def\PYZrb{]}
\makeatother


    % For linebreaks inside Verbatim environment from package fancyvrb.
    \makeatletter
        \newbox\Wrappedcontinuationbox
        \newbox\Wrappedvisiblespacebox
        \newcommand*\Wrappedvisiblespace {\textcolor{red}{\textvisiblespace}}
        \newcommand*\Wrappedcontinuationsymbol {\textcolor{red}{\llap{\tiny$\m@th\hookrightarrow$}}}
        \newcommand*\Wrappedcontinuationindent {3ex }
        \newcommand*\Wrappedafterbreak {\kern\Wrappedcontinuationindent\copy\Wrappedcontinuationbox}
        % Take advantage of the already applied Pygments mark-up to insert
        % potential linebreaks for TeX processing.
        %        {, <, #, %, $, ' and ": go to next line.
        %        _, }, ^, &, >, - and ~: stay at end of broken line.
        % Use of \textquotesingle for straight quote.
        \newcommand*\Wrappedbreaksatspecials {%
            \def\PYGZus{\discretionary{\char`\_}{\Wrappedafterbreak}{\char`\_}}%
            \def\PYGZob{\discretionary{}{\Wrappedafterbreak\char`\{}{\char`\{}}%
            \def\PYGZcb{\discretionary{\char`\}}{\Wrappedafterbreak}{\char`\}}}%
            \def\PYGZca{\discretionary{\char`\^}{\Wrappedafterbreak}{\char`\^}}%
            \def\PYGZam{\discretionary{\char`\&}{\Wrappedafterbreak}{\char`\&}}%
            \def\PYGZlt{\discretionary{}{\Wrappedafterbreak\char`\<}{\char`\<}}%
            \def\PYGZgt{\discretionary{\char`\>}{\Wrappedafterbreak}{\char`\>}}%
            \def\PYGZsh{\discretionary{}{\Wrappedafterbreak\char`\#}{\char`\#}}%
            \def\PYGZpc{\discretionary{}{\Wrappedafterbreak\char`\%}{\char`\%}}%
            \def\PYGZdl{\discretionary{}{\Wrappedafterbreak\char`\$}{\char`\$}}%
            \def\PYGZhy{\discretionary{\char`\-}{\Wrappedafterbreak}{\char`\-}}%
            \def\PYGZsq{\discretionary{}{\Wrappedafterbreak\textquotesingle}{\textquotesingle}}%
            \def\PYGZdq{\discretionary{}{\Wrappedafterbreak\char`\"}{\char`\"}}%
            \def\PYGZti{\discretionary{\char`\~}{\Wrappedafterbreak}{\char`\~}}%
        }
        % Some characters . , ; ? ! / are not pygmentized.
        % This macro makes them "active" and they will insert potential linebreaks
        \newcommand*\Wrappedbreaksatpunct {%
            \lccode`\~`\.\lowercase{\def~}{\discretionary{\hbox{\char`\.}}{\Wrappedafterbreak}{\hbox{\char`\.}}}%
            \lccode`\~`\,\lowercase{\def~}{\discretionary{\hbox{\char`\,}}{\Wrappedafterbreak}{\hbox{\char`\,}}}%
            \lccode`\~`\;\lowercase{\def~}{\discretionary{\hbox{\char`\;}}{\Wrappedafterbreak}{\hbox{\char`\;}}}%
            \lccode`\~`\:\lowercase{\def~}{\discretionary{\hbox{\char`\:}}{\Wrappedafterbreak}{\hbox{\char`\:}}}%
            \lccode`\~`\?\lowercase{\def~}{\discretionary{\hbox{\char`\?}}{\Wrappedafterbreak}{\hbox{\char`\?}}}%
            \lccode`\~`\!\lowercase{\def~}{\discretionary{\hbox{\char`\!}}{\Wrappedafterbreak}{\hbox{\char`\!}}}%
            \lccode`\~`\/\lowercase{\def~}{\discretionary{\hbox{\char`\/}}{\Wrappedafterbreak}{\hbox{\char`\/}}}%
            \catcode`\.\active
            \catcode`\,\active
            \catcode`\;\active
            \catcode`\:\active
            \catcode`\?\active
            \catcode`\!\active
            \catcode`\/\active
            \lccode`\~`\~
        }
    \makeatother

    \let\OriginalVerbatim=\Verbatim
    \makeatletter
    \renewcommand{\Verbatim}[1][1]{%
        %\parskip\z@skip
        \sbox\Wrappedcontinuationbox {\Wrappedcontinuationsymbol}%
        \sbox\Wrappedvisiblespacebox {\FV@SetupFont\Wrappedvisiblespace}%
        \def\FancyVerbFormatLine ##1{\hsize\linewidth
            \vtop{\raggedright\hyphenpenalty\z@\exhyphenpenalty\z@
                \doublehyphendemerits\z@\finalhyphendemerits\z@
                \strut ##1\strut}%
        }%
        % If the linebreak is at a space, the latter will be displayed as visible
        % space at end of first line, and a continuation symbol starts next line.
        % Stretch/shrink are however usually zero for typewriter font.
        \def\FV@Space {%
            \nobreak\hskip\z@ plus\fontdimen3\font minus\fontdimen4\font
            \discretionary{\copy\Wrappedvisiblespacebox}{\Wrappedafterbreak}
            {\kern\fontdimen2\font}%
        }%

        % Allow breaks at special characters using \PYG... macros.
        \Wrappedbreaksatspecials
        % Breaks at punctuation characters . , ; ? ! and / need catcode=\active
        \OriginalVerbatim[#1,codes*=\Wrappedbreaksatpunct]%
    }
    \makeatother

    % Exact colors from NB
    \definecolor{incolor}{HTML}{303F9F}
    \definecolor{outcolor}{HTML}{D84315}
    \definecolor{cellborder}{HTML}{CFCFCF}
    \definecolor{cellbackground}{HTML}{F7F7F7}

    % prompt
    \makeatletter
    \newcommand{\boxspacing}{\kern\kvtcb@left@rule\kern\kvtcb@boxsep}
    \makeatother
    \newcommand{\prompt}[4]{
        {\ttfamily\llap{{\color{#2}[#3]:\hspace{3pt}#4}}\vspace{-\baselineskip}}
    }
    

    
    % Prevent overflowing lines due to hard-to-break entities
    \sloppy
    % Setup hyperref package
    \hypersetup{
      breaklinks=true,  % so long urls are correctly broken across lines
      colorlinks=true,
      urlcolor=urlcolor,
      linkcolor=linkcolor,
      citecolor=citecolor,
      }
    % Slightly bigger margins than the latex defaults
    
    \geometry{verbose,tmargin=1in,bmargin=1in,lmargin=1in,rmargin=1in}
    
    

\begin{document}
    
    \maketitle
    
    

    
    \hypertarget{external-validation-clam-mb-uni2-h-pam50-pipeline}{%
\section{External Validation: CLAM-MB + UNI2-h PAM50
Pipeline}\label{external-validation-clam-mb-uni2-h-pam50-pipeline}}

This report evaluates the \textbf{CLAM-MB + UNI2-h} pipeline on two
external cohorts, \textbf{without retraining or domain adaptation}:

\begin{itemize}
\tightlist
\item
  \textbf{HSI-BC} -- 47 patients, one H\&E \texttt{.mrxs} WSI per
  patient, IHC-derived subtype labels.
\item
  \textbf{NOU} -- 28 patients, 203 tissue-block H\&E crops total,
  IHC-derived subtype labels.
\end{itemize}

The 10 CLAM-MB fold-checkpoints from TCGA-BRCA are reused as-is.
Ensemble predictions average softmax probabilities across all 10 folds.

\textbf{PAM50 classes:} LumA (0), LumB (1), Basal (2), Her2 (3).

Both cohorts use \textbf{IHC surrogate} annotations, which introduces
\textasciitilde{}10--20\% label noise vs.~genomic PAM50.

    \hypertarget{pipeline-overview}{%
\subsection{Pipeline Overview}\label{pipeline-overview}}

The pipeline is identical to the one in \texttt{evaluate\_pam50.pdf}:

\begin{enumerate}
\def\labelenumi{\arabic{enumi}.}
\tightlist
\item
  \textbf{Data preparation} -- IHC surrogate labels mapped to PAM50
  classes.
\item
  \textbf{Tissue patching} -- CLAM \texttt{create\_patches\_fp.py}.
  HSI-BC: 256px patches from \texttt{.mrxs}; NOU: 512px patches from
  TIFF crops.
\item
  \textbf{Feature extraction} -- UNI2-h (ViT-Giant/14, 1536-dim
  embeddings).
\item
  \textbf{Ensemble inference} -- 10 TCGA-BRCA CLAM-MB checkpoints,
  mean-softmax ensemble. No fine-tuning or domain adaptation.
\end{enumerate}

    \hypertarget{experimental-setup}{%
\subsection{Experimental setup}\label{experimental-setup}}

    \begin{Verbatim}[commandchars=\\\{\}]
[HSI-BC] loaded 47 ensemble predictions from 47 patients
[NOU] loaded 203 ensemble predictions from 28 patients
    \end{Verbatim}

    \hypertarget{cohort-characteristics}{%
\subsubsection{Cohort Characteristics}\label{cohort-characteristics}}

    \begin{Verbatim}[commandchars=\\\{\}]
Table 1: Cohort composition by PAM50 class
================================================================================
====================
                          Cohort  Slides Patients      LumA      LumB     Basal
Her2
TCGA-BRCA (training, 10-fold CV)     986     \textasciitilde{}930 508 (52\%) 222 (23\%) 173 (18\%)
83 (8\%)
               HSI-BC (external)      47       47  10 (21\%)  29 (62\%)    4 (9\%)
4 (9\%)
                  NOU (external)     203       28 152 (75\%)  23 (11\%)    6 (3\%)
22 (11\%)
================================================================================
====================
    \end{Verbatim}

    \hypertarget{per-fold-ensemble-member-performance}{%
\subsection{Per-Fold Ensemble Member
Performance}\label{per-fold-ensemble-member-performance}}

Each of the 10 TCGA-trained checkpoints is applied independently to
every external slide (no fold-wise splitting on the external set).

    \begin{Verbatim}[commandchars=\\\{\}]
Table 2: Per-fold ensemble member metrics on each external cohort
======================================================================
--- HSI-BC ---
        AUC    Acc   BAcc
Fold
0    0.7938 0.3617 0.4978
1    0.7708 0.2553 0.4220
2    0.7788 0.3830 0.5065
3    0.7809 0.2979 0.4345
4    0.7976 0.3617 0.5142
5    0.8196 0.4468 0.5159
6    0.7705 0.5957 0.5927
7    0.7778 0.4681 0.4871
8    0.7865 0.2979 0.4884
9    0.8050 0.4043 0.5315
  Mean +/- Std: AUC=0.7881 +/- 0.0158  Acc=0.3872 +/- 0.0992  BAcc=0.4991 +/-
0.0480

--- NOU ---
        AUC    Acc   BAcc
Fold
0    0.5121 0.4877 0.2192
1    0.5160 0.4089 0.2118
2    0.5269 0.5517 0.2405
3    0.4953 0.4926 0.2019
4    0.4742 0.4039 0.2097
5    0.4716 0.5961 0.2466
6    0.5141 0.4138 0.2037
7    0.4917 0.5369 0.2264
8    0.4772 0.5320 0.1966
9    0.5088 0.5665 0.2358
  Mean +/- Std: AUC=0.4988 +/- 0.0196  Acc=0.4990 +/- 0.0698  BAcc=0.2192 +/-
0.0174

======================================================================
    \end{Verbatim}

    \begin{center}
    \adjustimage{max size={0.9\linewidth}{0.9\paperheight}}{evaluate_external_files/evaluate_external_9_0.png}
    \end{center}
    { \hspace*{\fill} \\}
    
    \hypertarget{aggregated-ensemble-predictions-slide-level}{%
\subsection{Aggregated Ensemble Predictions
(Slide-Level)}\label{aggregated-ensemble-predictions-slide-level}}

Mean-softmax ensemble of 10 fold-checkpoints. For HSI-BC, slide-level =
patient-level (one WSI per patient). NOU patient-level aggregation
follows below.

    \begin{Verbatim}[commandchars=\\\{\}]
Table 3: Aggregated ensemble metrics (slide-level)
==============================================================================
  Cohort       Macro AUC       Acc      BAcc    Macro F1    Wt. F1
  ----------------------------------------------------------------
  HSI-BC          0.8042    0.4043    0.5315      0.3905    0.4004
  NOU             0.5028    0.5222    0.2215      0.2250    0.5600
==============================================================================

Note: BAcc (balanced accuracy) is the macro-averaged recall and equals
random-chance level (1/4 = 0.25) when the model collapses to a single class.
    \end{Verbatim}

    \begin{Verbatim}[commandchars=\\\{\}]
Classification report -- HSI-BC (slide-level)
============================================================
              precision    recall  f1-score   support

        LumA       0.33      0.60      0.43        10
        LumB       0.73      0.28      0.40        29
       Basal       0.50      0.25      0.33         4
        Her2       0.25      1.00      0.40         4

    accuracy                           0.40        47
   macro avg       0.45      0.53      0.39        47
weighted avg       0.58      0.40      0.40        47


Classification report -- NOU (slide-level)
============================================================
              precision    recall  f1-score   support

        LumA       0.79      0.66      0.72       152
        LumB       0.08      0.13      0.10        23
       Basal       0.00      0.00      0.00         6
        Her2       0.07      0.09      0.08        22

    accuracy                           0.52       203
   macro avg       0.23      0.22      0.23       203
weighted avg       0.61      0.52      0.56       203


    \end{Verbatim}

    \hypertarget{confusion-matrices}{%
\subsubsection{Confusion matrices}\label{confusion-matrices}}

    \begin{center}
    \adjustimage{max size={0.9\linewidth}{0.9\paperheight}}{evaluate_external_files/evaluate_external_14_0.png}
    \end{center}
    { \hspace*{\fill} \\}
    
    \hypertarget{roc-curves-one-vs-rest}{%
\subsubsection{ROC Curves (One-vs-Rest)}\label{roc-curves-one-vs-rest}}

    \begin{center}
    \adjustimage{max size={0.9\linewidth}{0.9\paperheight}}{evaluate_external_files/evaluate_external_16_0.png}
    \end{center}
    { \hspace*{\fill} \\}
    
    \hypertarget{per-class-breakdown}{%
\subsubsection{Per-class breakdown}\label{per-class-breakdown}}

    \begin{Verbatim}[commandchars=\\\{\}]
Table 4: Per-class metrics on external cohorts (slide-level)
==============================================================================
Cohort Class  Support    AUC  Precision  Recall     F1
HSI-BC  LumA       10 0.6486     0.3333  0.6000 0.4286
HSI-BC  LumB       29 0.6552     0.7273  0.2759 0.4000
HSI-BC Basal        4 0.9651     0.5000  0.2500 0.3333
HSI-BC  Her2        4 0.9477     0.2500  1.0000 0.4000
   NOU  LumA      152 0.5702     0.7891  0.6645 0.7214
   NOU  LumB       23 0.4529     0.0833  0.1304 0.1017
   NOU Basal        6 0.4272     0.0000  0.0000 0.0000
   NOU  Her2       22 0.5608     0.0667  0.0909 0.0769
==============================================================================
    \end{Verbatim}

    \hypertarget{nou-patient-level-evaluation}{%
\subsection{NOU Patient-Level
Evaluation}\label{nou-patient-level-evaluation}}

We aggregate the ensemble softmax probabilities across all tissue blocks
of a patient by averaging, then argmax for the patient-level prediction.

    \begin{Verbatim}[commandchars=\\\{\}]
NOU patient-level evaluation: n = 28 patients, median 7 blocks/patient
============================================================
  Macro AUC:           0.4246
  Accuracy:            0.7143
  Balanced accuracy:   0.2381
  Macro F1:            0.2128
  Weighted F1:         0.6383

              precision    recall  f1-score   support

        LumA       0.77      0.95      0.85        21
        LumB       0.00      0.00      0.00         3
       Basal       0.00      0.00      0.00         1
        Her2       0.00      0.00      0.00         3

    accuracy                           0.71        28
   macro avg       0.19      0.24      0.21        28
weighted avg       0.58      0.71      0.64        28

    \end{Verbatim}

    \begin{center}
    \adjustimage{max size={0.9\linewidth}{0.9\paperheight}}{evaluate_external_files/evaluate_external_21_0.png}
    \end{center}
    { \hspace*{\fill} \\}
    
    \hypertarget{comparison-internal-vs.-external}{%
\subsection{Comparison: Internal
vs.~External}\label{comparison-internal-vs.-external}}

    \begin{Verbatim}[commandchars=\\\{\}]
Table 5: TCGA-BRCA internal vs. HSI-BC and NOU external
================================================================================
==========
               TCGA-BRCA (10-fold CV)  HSI-BC (slide=patient)  NOU (slide-level)
NOU (patient-level)
Macro AUC                      0.8879                  0.8042             0.5028
0.4246
Accuracy                       0.7181                  0.4043             0.5222
0.7143
Balanced Acc.                     NaN                  0.5315             0.2215
0.2381
Weighted F1                    0.7179                  0.4004             0.5600
0.6383
LumA AUC                       0.8967                  0.6486             0.5702
0.6599
LumB AUC                       0.8174                  0.6552             0.4529
0.4667
Basal AUC                      0.9529                  0.9651             0.4272
0.1852
Her2 AUC                       0.8846                  0.9477             0.5608
0.3867
================================================================================
==========

TCGA values are reproduced from evaluate\_pam50.pdf (10-fold CV, aggregated test
set).
    \end{Verbatim}

    \hypertarget{findings}{%
\subsection{Findings}\label{findings}}

\begin{enumerate}
\def\labelenumi{\arabic{enumi}.}
\item
  \textbf{HSI-BC: useful transfer, substantial drop.} Macro AUC drops
  from 0.888 (TCGA internal) to \textbf{0.804} (\textasciitilde{}9 pp).
  Per-class AUC ranking preserved (Basal, Her2, LumA, LumB). Low
  accuracy (0.40) reflects LumB dominance (62\% of cohort) -- the model
  frequently confuses LumB with LumA.
\item
  \textbf{NOU: model collapses to chance.} Slide-level macro AUC = 0.50,
  patient-level = 0.42. The 71\% patient-level accuracy is misleading --
  it merely reflects LumA prevalence (75\%). Balanced accuracy is
  \textasciitilde{}0.24, i.e.~random. All 10 ensemble members produce
  AUC in {[}0.47, 0.53{]}, confirming systematic failure.
\item
  \textbf{NOU collapse is likely a domain shift issue.} NOU TIFFs are
  small tissue crops (\textasciitilde{}210 patches/slide
  vs.~\textasciitilde{}48k for HSI-BC). CLAM attention, trained on large
  bags, cannot extract signal from bags two orders of magnitude smaller.
\item
  \textbf{Class imbalance amplifies AUC--accuracy gap.} Both cohorts are
  heavily luminal-skewed (HSI-BC 83\%, NOU 89\%), with 4 or fewer
  Basal/Her2 patients each. Per-class metrics have wide uncertainty at
  these sample sizes.
\item
  \textbf{IHC-to-PAM50 label noise caps performance.} Both cohorts use
  IHC surrogates (\textasciitilde{}80\% concordance with genomic PAM50),
  so some ``errors'' are label-definition disagreements.
\end{enumerate}

    \hypertarget{next-steps}{%
\subsection{Next Steps}\label{next-steps}}

\begin{enumerate}
\def\labelenumi{\arabic{enumi}.}
\tightlist
\item
  \textbf{Re-patch NOU} with a patching strategy that matches the TCGA
  training distribution (larger bags, proper field-of-view alignment).
\item
  \textbf{Stain normalisation} (Macenko/Vahadane) before feature
  extraction to reduce domain shift.
\item
  \textbf{Bootstrap confidence intervals} (stratified, 1000 resamples)
  for all headline metrics.
\item
  \textbf{Light fine-tuning} of CLAM-MB on external data to distinguish
  feature-level vs.~aggregator-level failure.
\item
  \textbf{Multimodal fusion} using clinico-pathological features (HSI-BC
  clinical data, NOU CTC EP/EMT counts) alongside WSI embeddings.
\end{enumerate}

    

    \hypertarget{summary}{%
\subsection{Summary}\label{summary}}

    \begin{Verbatim}[commandchars=\\\{\}]
======================================================================
  CLAM-MB + UNI2-h | PAM50 Subtyping | External Validation
  Model: 10x checkpoints from TCGA-BRCA (no retraining)
======================================================================

--- TCGA-BRCA (internal benchmark from evaluate\_pam50.pdf) ---
  Macro AUC       0.8879
  Accuracy        0.7181
  Weighted F1     0.7179
  LumA AUC        0.8967
  LumB AUC        0.8174
  Basal AUC       0.9529
  Her2 AUC        0.8846

--- HSI-BC (slide-level, n = 47) ---
  Macro AUC      0.8042
  Accuracy       0.4043
  Balanced Acc.  0.5315
  Weighted F1    0.4004

--- NOU (slide-level, n = 203) ---
  Macro AUC      0.5028
  Accuracy       0.5222
  Balanced Acc.  0.2215
  Weighted F1    0.5600

--- NOU (patient-level, n = 28) ---
  Macro AUC      0.4246
  Accuracy       0.7143
  Balanced Acc.  0.2381
  Weighted F1    0.6383
======================================================================
    \end{Verbatim}


    % Add a bibliography block to the postdoc
    
    
    
\end{document}
